\chapter{Implementación}
\label{cap:implementacion}

Este capítulo describe cómo se ha implementado, en la práctica, el diseño
presentado en el capítulo~\ref{cap:diseno}. El capítulo se cierra con una sección
dedicada explícitamente a los enfoques que se exploraron y finalmente no se
incorporaron al sistema.

\section{Entorno y herramientas}
\label{sec:implementacion-entorno}

Todo el desarrollo se ha realizado en Python 3.11, dentro de un entorno
virtual (\texttt{venv\_tfg}) situado dentro del propio repositorio, con las
dependencias fijadas en
\texttt{requirements.txt} para garantizar la reproducibilidad.
La Tabla~\ref{tab:impl-librerias} recoge las librerías con un papel directo
en el sistema. Se omiten las dependencias transitivas de bajo nivel, como
\texttt{torch} o \texttt{pandas}, que \texttt{requirements.txt} incluye por
ser un volcado completo del entorno pero que ningún
script del proyecto importa directamente: llegan al entorno como
dependencias de \texttt{boxmot} y de \texttt{ultralytics}.

\clearpage

\begin{table}[htp]
\centering
\caption{Librerías principales empleadas y su función en el sistema.}
\label{tab:impl-librerias}
\renewcommand{\arraystretch}{1.25}
\begin{tabular}{|p{2.6cm}|p{2.5cm}|p{7.0cm}|}
\hline
\rowcolor{gray!25}
\textbf{Librería} & \textbf{Versión} & \textbf{Función en el sistema} \\
\hline
\texttt{boxmot} & 10.0.43 & Implementaciones de ByteTrack y OC-SORT bajo una
  interfaz común. Versión fijada
  deliberadamente por el problema P4 de la
  Tabla~\ref{tab:impl-problemas}. \\
\hline
\texttt{trackeval} & \texttt{sn-trackeval} 0.4.0 & Evaluación con
  métricas MOT estándar
  sección. Fork de TrackEval adaptado por
  SoccerNet al formato de su benchmark. \\
\hline
\texttt{opencv-python} & 4.13.0.92 & Lectura y anotación de frames,
  cálculo de la homografía, (sección~\ref{sec:implementacion-homografia}). \\
\hline
\texttt{filterpy} & 1.4.5 & Implementación del filtro de Kalman
  (sección~\ref{subsec:fundamentos-kalman}) usada internamente por los
  trackers de \texttt{boxmot}. \\
\hline
\texttt{lapx} & 0.9.4 & Resolución del problema de asignación mediante el
  algoritmo húngaro (sección~\ref{subsec:fundamentos-hungaro}), usada
  internamente por los trackers de \texttt{boxmot}. \\
\hline
\texttt{scikit-learn} & 1.8.0 & K-means para la clasificación de equipos
  (sección~\ref{sec:fundamentos-kmeans}). \\
\hline
\texttt{scipy} & 1.17.1 & Estimación de densidad por núcleos
  (\texttt{gaussian\_kde}) en los mapas de calor. \\
\hline
\texttt{matplotlib} & 3.10.9 & Generación de todas las figuras y gráficas
  del proyecto. \\
\hline
\end{tabular}
\end{table}

Cabe señalar una inconsistencia menor detectada al redactar esta memoria:
aunque \texttt{trackeval} está instalado y en uso real en el entorno de
desarrollo (a través del paquete \texttt{sn-trackeval}), no aparece listado
de forma explícita en \texttt{requirements.txt}, lo que sugiere que ese
fichero no se regeneró tras instalarlo por separado. No afecta al
funcionamiento del proyecto en el equipo de desarrollo actual, pero es una
brecha de reproducibilidad frente al requisito RNF-001 que convendría
corregir --regenerando \texttt{requirements.txt} con \texttt{sn-trackeval}
explícitamente incluido-- antes de que un tercero intente reproducir el
entorno desde cero.

\subsection{Traslados de equipo de desarrollo}

El desarrollo no se ha mantenido en un único equipo, lo que ha exigido
recrear el entorno de trabajo en dos ocasiones: el proyecto se inició en un
portátil personal, se trasladó a mitad de desarrollo a un ordenador de
sobremesa con GPU AMD Radeon RX 6900 XT (16\,GB VRAM) en Windows, y finalmente a un equipo macOS. Cada
traslado se apoyó en el mismo mecanismo: dependencias fijadas en
\texttt{requirements.txt} y código versionado en Git, lo que permitió
recrear el entorno funcional en el nuevo equipo sin más incidencias que las
propias de cada sistema operativo.

\subsection{Problemas de entorno encontrados}
\label{subsec:implementacion-problemas}

Configurar y mantener este entorno no fue inmediato. La
Tabla~\ref{tab:impl-problemas} documenta, siguiendo el
mismo criterio de transparencia del resto de esta memoria, los catorce problemas de
entorno registrados durante el desarrollo, con su solución.

\begin{table}[htp]
\centering
\caption{Problemas de entorno encontrados durante el desarrollo y su solución.}
\label{tab:impl-problemas}
\renewcommand{\arraystretch}{1.05}
\footnotesize
\begin{tabular}{|p{0.7cm}|p{5.5cm}|p{5.3cm}|}
\hline
\rowcolor{gray!25}
\textbf{ID} & \textbf{Problema} & \textbf{Solución} \\
\hline
P1 & Python 3.14 incompatible con las librerías del proyecto &
  Fijar Python 3.11 en el entorno virtual \\
\hline
P2 & Mezcla de entornos conda y venv activos a la vez, con conflictos de
  paquetes & Desactivar conda y trabajar exclusivamente con \texttt{venv} \\
\hline
P3 & \texttt{conda activate venv\_tfg} fallaba, ya que \texttt{venv\_tfg}
  no es un entorno conda & Activar el entorno con el mecanismo propio de
  \texttt{venv} (\texttt{source venv\_tfg/bin/activate}) \\
\hline
P4 & \texttt{boxmot} v18 no expone las clases de cada tracker por separado,
  necesarias para la comparativa & Fijar la versión \texttt{10.0.43} en
  \texttt{requirements.txt} \\
\hline
P5 & Error de tipo en \texttt{kalman\_filter.py} de OC-SORT por
  incompatibilidad con la versión de NumPy instalada & Fijar
  \texttt{numpy==1.26.4} en el momento de detectarlo\footnotemark \\
\hline
P6 & Detecciones de \texttt{det.txt} con 5 columnas en lugar de las 6 que
  espera \texttt{boxmot} & Añadir manualmente la columna de clase (valor
  \texttt{0}) antes de pasar las detecciones al tracker \\
\hline
P7 & La variable de entorno \texttt{\$HOME} se colaba en comandos al
  copiar y pegar instrucciones entre terminales & Revisar manualmente los
  comandos antes de ejecutarlos \\
\hline
P8 & El entorno virtual se creó fuera de la carpeta del proyecto y VS Code
  no lo detectaba automáticamente & Recrear el entorno dentro del
  repositorio \\
\hline
P9 & Los entornos virtuales de Python no son reubicables (las rutas
  absolutas quedan grabadas en sus scripts de activación) &
  Recrearlos siempre en su ubicación final, no moverlos tras crearlos \\
\hline
P10 & \texttt{requirements.txt} estaba vacío & Regenerarlo con
  \texttt{pip freeze > requirements.txt} \\
\hline
P13 & \texttt{load\_dotenv()} no recargaba el fichero \texttt{.env} tras
  modificarlo & Forzar la recarga con \texttt{load\_dotenv(override=True)} \\
\hline
P14 & Carpetas basura generadas por rutas mal resueltas
  (\texttt{path/}, \texttt{\$HOME/}, \texttt{.venv/}, \texttt{models/}) &
  Identificarlas y eliminarlas manualmente \\
\hline
P15 & Traslado del proyecto a macOS: posible incompatibilidad de entorno &
  Entorno recreado desde \texttt{requirements.txt} sin incidencias
  destacables; solo hubo que ajustar alguna ruta de Windows a macOS \\
\hline
P17 & Los vídeos exportados con el \emph{fourcc} \texttt{mp4v} de OpenCV no
  se reproducían en el navegador & Usar el \emph{fourcc} \texttt{avc1}
  (H.264), disponible porque OpenCV, con esto, el tamaño de cada vídeo se redujo en torno a un 75\,\% \\
\hline
\end{tabular}
\end{table}
\footnotetext{El \texttt{requirements.txt} final del proyecto fija
\texttt{numpy==2.4.4}, una versión posterior a la que resolvió este
problema en su momento. Esto sugiere que la incompatibilidad original fue
resuelta más adelante en versiones posteriores de \texttt{boxmot} o de
OC-SORT, aunque no se ha vuelto a verificar explícitamente la causa
original tras la actualización.}

De estos catorce problemas, la mayoría son propios de la configuración
inicial de cualquier entorno Python (P1-P3, P8-P10) y no volvieron a
aparecer una vez resueltos. Los más relevantes para entender decisiones
posteriores de este proyecto son P4 y P6.
P17, ya en la fase de aplicación web, muestra que ese mismo tipo de
fricción de entorno no es exclusivo de Python: elegir un códec de vídeo
compatible con el navegador es un problema de la misma naturaleza en un
\emph{stack} tecnológico distinto.

\section{Estructura del repositorio}
\label{sec:implementacion-estructura}

El código del pipeline offline reside en \texttt{src/}, organizado en
paquetes por responsabilidad, resumidos en la
Tabla~\ref{tab:impl-estructura}.

\begin{table}[htp]
\centering
\caption{Paquetes del repositorio y su responsabilidad.}
\label{tab:impl-estructura}
\renewcommand{\arraystretch}{1.25}
\begin{tabular}{|p{3.6cm}|p{8.3cm}|}
\hline
\rowcolor{gray!25}
\textbf{Paquete} & \textbf{Responsabilidad} \\
\hline
\texttt{src/\allowbreak visualizacion/} & Scripts exploratorios iniciales sobre el
  \emph{ground truth}. \\
\hline
\texttt{src/\allowbreak tracking/} & Implementación de ByteTrack y OC-SORT, ejecución
  batch sobre un split completo, y exportación de vídeos anotados
  (\texttt{src/\allowbreak tracking/\allowbreak export/}). \\
\hline
\texttt{src/\allowbreak evaluation/} & Generación de \emph{seqmaps} y evaluación con
  \texttt{TrackEval} sobre train y test. \\
\hline
\texttt{src/\allowbreak hyperparameter\_\allowbreak study/} & Barrido de hiperparámetros de cada
  algoritmo, configuración óptima y resumen comparativo. \\
\hline
\texttt{src/\allowbreak metadata/} & Todo lo relacionado con la generación y gestión de metadatos. \\
\hline
\texttt{src/\allowbreak utils/} & Scripts de utilidad para tareas de un solo uso. \\
\hline
\texttt{src/\allowbreak web\_\allowbreak export/} & Generación de metadatos de OC-SORT y
  empaquetado de todos los artefactos estáticos que consume la aplicación
  web (sección~\ref{sec:implementacion-web}). \\
\hline
\end{tabular}
\end{table}

El código de la aplicación web, al no ser Python, vive fuera de
\texttt{src/}: en \texttt{frontend/}, un proyecto independiente de React.(sección~\ref{sec:implementacion-web}).

\section{Implementación del tracking}
\label{sec:implementacion-tracking}

\subsection{Scripts de visualización inicial}
\label{subsec:impl-visualizacion}

Antes de implementar ningún algoritmo de tracking, la primera fase del
proyecto se dedicó a comprender el
formato del dataset mediante scripts exploratorios, todos
ellos trabajando directamente sobre el \emph{ground truth}
(\texttt{gt.txt}) y no sobre predicciones reales de ningún tracker:
\texttt{reproducir\_vid.py} reproduce una secuencia como vídeo con OpenCV;
\texttt{bboxes\_ids.py} dibuja las cajas del \emph{ground truth} en verde
con su identificador; y \texttt{bboxes\_ids\_dif\_colors.py}
añade un color único y persistente por identificador, para facilitar el
seguimiento visual de cada jugador entre frames.

El último de estos scripts exploratorios,
\texttt{bboxes\_balltrack\_geoprop.py}, ya incorpora el filtro geométrico
de separación del balón (sección~\ref{subsec:impl-balon}) y dibuja el balón
en rojo, distinto de los jugadores. Es interesante notar que
\texttt{bboxes\_ids.py}.

\subsection{Filtrado geométrico del balón}
\label{subsec:impl-balon}

Como se diseñó en la sección~\ref{subsec:diseno-balon}, el balón se separa
de las detecciones de personas mediante dos condiciones sobre su caja
delimitadora, aplicadas de forma idéntica en todos los scripts del
proyecto que necesitan distinguir balón de jugadores:

\begin{center}
\texttt{BALL\_MAX\_AREA = 2000} \quad
\texttt{BALL\_MIN\_RATIO = 0.7} \quad
\texttt{BALL\_MAX\_RATIO = 1.3}
\end{center}

Una detección se clasifica como balón cuando su área en píxeles
($w \times h$) es inferior a \texttt{BALL\_MAX\_AREA} \emph{y}, a la vez, su
proporción $w/h$ está dentro del rango
[\texttt{BALL\_MIN\_RATIO}, \texttt{BALL\_MAX\_RATIO}], es decir, cuando la
caja es aproximadamente cuadrada. Estos tres valores se fijaron de forma
que, observando las dimensiones típicas del balón frente a las de un
jugador en las secuencias de SoccerNet, y se mantienen constantes en todo
el proyecto para que el balón se identifique de forma consistente en todas
las fases.

\subsection{ByteTrack y OC-SORT}
\label{subsec:impl-algoritmos}

Ambos algoritmos se implementan siguiendo el bucle diseñado en la
sección~\ref{subsec:diseno-tracking-algoritmos}, instanciando en cada caso la clase
correspondiente de \texttt{boxmot}: \texttt{BYTETracker()} en
\texttt{bytetracker\_alg.py}, con su configuración por defecto como punto
de partida, y
\texttt{OCSORT} también con su configuración por defecto en
\texttt{sort\_alg.py}.

El nombre \texttt{sort\_alg.py}, es una marca, documentada en la sección~\ref{subsec:fundamentos-sort}. La
intención original era implementar el SORT clásico, pero \texttt{boxmot}
en su versión fijada no expone una clase \texttt{SORT} independiente,
solo su evolución directa \texttt{OCSORT}. El fichero conserva su nombre
original como constancia de esa decisión, aunque su contenido implementa
realmente OC-SORT.

Las detecciones que llegan de \texttt{det.txt} en formato
$(x, y, w, h)$ se convierten al formato $(x_1, y_1, x_2, y_2, \text{conf},
\text{clase})$ que espera \texttt{boxmot} antes de invocar
\texttt{tracker.update()} en cada frame, y el resultado se reasigna a una numeración secuencial
mediante un diccionario de correspondencia, antes de escribirse en
el fichero de resultados con el formato MOT estándar.

\subsection{Ejecución batch y reorganización de resultados}

La ejecución sobre la totalidad de un split se implementa en
\texttt{run\_all\_train.py} y \texttt{run\_all\_test.py}, que recorren
todas las secuencias del directorio del split correspondiente y aplican,
para cada una y cada algoritmo configurado, el mismo bucle de tracking
descrito en la sección anterior (Figura~\ref{fig:diseno-batch}). Ambos
scripts comparten la misma lógica de filtrado del balón y de escritura de
resultados que las versiones de un único algoritmo y una única secuencia.

En una primera versión del proyecto, los
resultados de cada secuencia se escribían en una subcarpeta propia
(\texttt{results/bytetracker/SNMOT-060/SNMOT-060.txt}), estructura que
\texttt{TrackEval} no esperaba de forma directa. Este script recorrió una
única vez los resultados ya generados y los aplanó a la estructura que
finalmente adopta el proyecto (\texttt{results/bytetracker/SNMOT-060.txt}),
sin necesidad de rehacer el tracking desde cero. No forma parte del flujo
habitual de ejecución del sistema, pero se conserva en el repositorio como
constancia de esa reorganización.

\section{Implementación de la evaluación}
\label{sec:implementacion-evaluacion}

La evaluación se
implementa en dos scripts,
\texttt{evaluate\_train.py} y \texttt{evaluate\_test.py}, que configuran
\texttt{TrackEval} con los mismos parámetros salvo el split evaluado:

\begin{center}
\begin{minipage}{0.85\textwidth}
\begin{verbatim}
dataset_config['BENCHMARK']          = 'SoccerNet'
dataset_config['TRACKERS_TO_EVAL']   = ['bytetracker', 'ocsort']
dataset_config['SKIP_SPLIT_FOL']     = True
dataset_config['DO_PREPROC']         = False
metrics_config = {'METRICS': ['HOTA', 'MOTA', 'IDF1'], 'THRESHOLD': 0.5}
\end{verbatim}
\end{minipage}
\end{center}

donde \texttt{BENCHMARK = 'SoccerNet'} indica a \texttt{TrackEval} que use
el \emph{seqmap} generado por
\texttt{seqmaps\_creation\_test.py}/\texttt{seqmpas\_creation\_train.py}
(sección~\ref{sec:diseno-evaluacion}) en lugar de un benchmark predefinido,
y \texttt{THRESHOLD = 0.5} fija el umbral de IoU al que se calculan MOTA e
IDF1, ambas métricas de umbral único.

\subsection{Aclaración metodológica: HOTA frente a HOTA(0)}
\label{subsec:impl-hota-vs-hota0}

Antes de presentar cualquier resultado numérico en el
capítulo~\ref{cap:pruebas}, es necesario dejar constancia de un matiz
metodológico detectado al preparar esta memoria, en cumplimiento del
requisito de documentación honesta. \texttt{TrackEval} exporta,
entre sus columnas de salida, tanto \textbf{HOTA}, promediada sobre el
rango completo de umbrales de IoU de 0,05 a 0,95, como
\textbf{HOTA(0)}\footnote{La notación \texttt{HOTA(0)} es la que usa
\texttt{TrackEval} internamente para referirse al valor de HOTA en el
primer umbral de su rango de evaluación, que con la configuración de este
proyecto (\texttt{THRESHOLD = 0.5}) coincide con el umbral de 0,5 usado
también por MOTA e IDF1.}.

El script de resumen del estudio de hiperparámetros
(\texttt{summarize\_hyperparam.py}) lee explícitamente la columna
\texttt{HOTA(0)} de cada resultado:

\begin{center}
\begin{minipage}{0.8\textwidth}
\begin{verbatim}
hota = float(row.get('HOTA(0)', row.get('HOTA___50', 0))) * 100
\end{verbatim}
\end{minipage}
\end{center}

lo que significa que \textbf{todos los valores de ``HOTA'' reportados en
el estudio de hiperparámetros de este proyecto}
\textbf{son en realidad HOTA(0)}, es decir, HOTA evaluada únicamente al
50\,\% de solapamiento IoU, y no la métrica HOTA integrada que
describe la sección~\ref{subsec:fundamentos-hota}. Ambas son métricas
válidas y ambas las calcula \texttt{TrackEval} de forma estándar, pero no
son intercambiables: al integrar sobre un rango de umbrales en lugar de
uno solo, la HOTA completa es más exigente y, en general, produce un valor
distinto.
El capítulo~\ref{cap:pruebas} presenta ambos valores por separado allí
donde la distinción afecta a las conclusiones, y este matiz se retoma
explícitamente donde corresponda para no inducir a una lectura equivocada
de los resultados.

\section{Implementación del estudio de hiperparámetros}
\label{sec:implementacion-hiperparametros}

El estudio de hiperparámetros se aplica siempre sobre el
conjunto de \textbf{test}, para que la
configuración resultante no esté ajustada a datos ya vistos durante el
desarrollo inicial. Cada algoritmo tiene su propia configuración base y su
propio rango de búsqueda por hiperparámetro, recogidos en la
Tabla~\ref{tab:impl-hiperparametros-espacio}.

\begin{table}[htp]
\centering
\caption{Configuración base y rango de búsqueda de cada hiperparámetro.}
\label{tab:impl-hiperparametros-espacio}
\renewcommand{\arraystretch}{1.25}
\begin{tabular}{|p{2.6cm}|p{3.1cm}|p{1.6cm}|p{5.7cm}|}
\hline
\rowcolor{gray!25}
\textbf{Algoritmo} & \textbf{Hiperparámetro} & \textbf{Base} & \textbf{Rango de búsqueda} \\
\hline
\multirow{3}{*}{ByteTrack} & \texttt{track\_thresh} & 0.45 & 0.2, 0.3, 0.45, 0.6 \\
\cline{2-4}
 & \texttt{match\_thresh} & 0.8 & 0.6, 0.7, 0.8, 0.9 \\
\cline{2-4}
 & \texttt{track\_buffer} & 25 & 10, 20, 25, 40 \\
\hline
\multirow{4}{*}{OC-SORT} & \texttt{det\_thresh} & 0.3 & 0.1, 0.2, 0.3, 0.4, 0.5 \\
\cline{2-4}
 & \texttt{max\_age} & 30 & 10, 20, 30, 50 \\
\cline{2-4}
 & \texttt{min\_hits} & 3 & 1, 2, 3, 5 \\
\cline{2-4}
 & \texttt{iou\_threshold} & 0.3 & 0.1, 0.2, 0.3, 0.4, 0.5 \\
\hline
\end{tabular}
\end{table}

Los hiperparámetros \texttt{track\_thresh} (ByteTrack) y
\texttt{det\_thresh} (OC-SORT) no mostraron ningún efecto sobre la
métrica en el rango probado. El análisis de sensibilidad completo del
resto de hiperparámetros, con el efecto observado de cada uno sobre
HOTA(0), se presenta en la sección~\ref{sec:pruebas-hiperparametros} del
capítulo~\ref{cap:pruebas}, junto con la interpretación de esos
resultados: esta sección se limita a describir cómo se generó y validó la
configuración óptima.

A partir de esos resultados individuales, \texttt{bytetracker\_optimal.py}
y \texttt{ocsort\_optimal.py} fijan y validan de forma independiente la
configuración combinada que RF-008 exige confirmar:
\texttt{match\_thresh=0.9, track\_buffer=10} para ByteTrack (manteniendo
\texttt{track\_thresh=0.45}, sin efecto), y
\texttt{iou\_threshold=0.1, max\_age=10, min\_hits=1} para OC-SORT
(manteniendo \texttt{det\_thresh=0.3}, sin efecto). Los resultados
numéricos de esta validación independiente, y la comparación entre ambas
configuraciones óptimas, se presentan igualmente en la
sección~\ref{sec:pruebas-hiperparametros}; esta configuración de
ByteTrack es, en cualquier caso, la que se usa como base de todo el
pipeline de metadatos de la sección~\ref{sec:implementacion-metadatos}.

Por último, \texttt{summarize\_hyperparam.py} consolida todos estos
resultados en \texttt{evaluation\_hyperparam/summary\_table.csv} y genera
una gráfica de barras por hiperparámetro y algoritmo en
\texttt{evaluation\_hyperparam/plots/}.

\section{Implementación del pipeline de metadatos}
\label{sec:implementacion-metadatos}

Todo el pipeline de metadatos parte de los resultados de tracking generados
con la configuración óptima de ByteTrack, como se
justificó en la sección anterior.

\subsection{Clasificación de equipos}
\label{subsec:impl-equipos}

El script \texttt{team\_classification.py} implementa el diseño en dos niveles de
K-means de la sección~\ref{subsec:diseno-equipos}. Para limitar el coste
computacional, no procesa todos los frames de la secuencia, sino un
muestreo de \texttt{1 de cada 5} (\texttt{SAMPLE\_EVERY = 5}); dado que un
jugador aparece en decenas de frames a lo largo de una secuencia, este
muestreo basta para determinar su equipo mayoritario sin analizar la
secuencia completa. Para cada detección muestreada, se recorta la región
superior de su caja delimitadora para aislar la
zona de la camiseta y evitar incluir piernas, césped o el balón en
el recorte, antes de aplicar K-means con $k=1$ para obtener su color
dominante en espacio de color BGR. Los colores dominantes de todas las
muestras de la secuencia se agrupan después con K-means $k=2$, y el equipo
final de cada \texttt{track\_id} se determina por mayoría entre sus
muestras individuales, así evita que el ruido puede clasificar mal alguna muestra
puntual sin afectar al resultado global.

\subsection{Mapas de calor}

El pipeline de mapas de calor tiene, en realidad, dos implementaciones, resultado de una iteración de diseño explícita
documentada aquí: \texttt{general\_heatmap.py},
la primera versión, representa las posiciones directamente en píxeles de
cámara mediante \texttt{hexbin} de \texttt{matplotlib}, sin proyectar sobre
ningún campo dibujado; \texttt{heatmap\_field.py}, la versión final y la
que se usa en el resto del pipeline, sustituye \texttt{hexbin}
por una estimación de densidad por núcleos (\texttt{gaussian\_kde} de
\texttt{scipy}) proyectada sobre un campo de 105\,m $\times$ 68\,m dibujado
a escala (Figura~\ref{fig:fundamentos-campo}), con dos mejoras
concretas sobre la primera versión: un degradado continuo que también
representa las zonas de baja densidad y un escalado por el rango
observado de posiciones (percentiles 5 y 95). El script \texttt{heatmap\_by\_team.py}
reutiliza exactamente esta misma lógica de proyección y KDE ya validada en
\texttt{heatmap\_field.py}, aplicándola por separado a las posiciones de
cada equipo pero con una única escala
común a ambos, para que los dos mapas resultantes sean directamente
comparables entre sí (Figura~\ref{fig:diseno-heatmap-ejemplo}).

\subsection{Homografía real: distancia y velocidad}
\label{sec:implementacion-homografia}

Como se diseñó en la sección~\ref{subsec:diseno-homografia}, la homografía
real se calcula sobre un tramo concreto de cámara estable de cada secuencia
de demostración, identificando manualmente 4 puntos
de referencia en un frame representativo
del tramo. Para SNMOT-116, el tramo usado es el intervalo de frames
$[1, 380]$, correspondiente a una jugada de córner con cámara estable; para
SNMOT-118, el intervalo $[350, 650]$. Los puntos de imagen se identifican
sobre un frame de referencia de $1456 \times 816$ píxeles y se reescalan a
la resolución real de las secuencias ($1920 \times 1080$) antes de calcular
la homografía, y se corresponden con las coordenadas reales del área
pequeña según las medidas FIFA.

Antes de usarse, la matriz de homografía calculada se valida
reproyectando estos mismos 4 puntos y comparando el resultado con las
coordenadas de campo esperadas: para el tramo de SNMOT-116, el error medio
de reproyección obtenido es de apenas unos pocos centímetros, lo que
confirma que la matriz describe correctamente la transformación en la zona
próxima a los puntos de calibración.
A partir de esa matriz, la posición de cada jugador se proyecta a
metros reales frame a frame, y de esa secuencia de posiciones proyectadas
se calcula la distancia total recorrida y la velocidad, tanto media como de pico.

Para los mapas de calor calibrados con esta homografía, el diseño incorpora además el
sombreado visual descrito en la
sección~\ref{subsec:diseno-homografia}
(Figura~\ref{fig:diseno-heatmap-homografia-ejemplo}): las zonas del campo
alejadas de los 4 puntos de calibración se oscurecen para indicar
visualmente que la proyección en esa región no es fiable. Para SNMOT-118,
además, dos identificadores de tracking
(\texttt{EXCLUDED\_TRACK\_IDS = [328, 274]}) se excluyen manualmente del
mapa de calor calibrado: corresponden a jugadores cuya posición proyectada
se aleja del campo de forma incompatible con un movimiento humano real,
señal de que su trayectoria de tracking está fuera de la
zona de cámara estable para la que se calculó la homografía.

\subsection{Posesión de balón}
\label{subsec:impl-posesion}

El script \texttt{ball\_possesion.py} implementa la heurística de proximidad diseñada
en la sección~\ref{subsec:diseno-posesion}: para cada frame en que el
filtro geométrico detecta un balón en el
\texttt{det.txt} original, se calcula la
distancia euclídea entre su centro y el centro de cada jugador presente en
ese frame según los resultados de tracking, y se asigna la posesión al
equipo del jugador más cercano siempre que esa distancia mínima no supere
\texttt{MAX\_POSSESSION\_DISTANCE\_PX = 150} píxeles; por encima
de ese umbral, el balón se considera suelto y el frame no computa para
ningún equipo. El resultado se agrega tanto a nivel de frame
(\texttt{\_possession.csv}) como agregado en porcentaje por equipo
(\texttt{\_possession\_summary.csv}).

\section{Implementación de la aplicación web}
\label{sec:implementacion-web}

Esta sección detalla la implementación de la aplicación web diseñada en la
sección~\ref{sec:diseno-web}: los dos scripts nuevos de
\texttt{src/web\_export/} que preparan los datos, y el proyecto
\texttt{frontend/} que los consume.

\subsection{Generación de metadatos para OC-SORT}
\label{subsec:impl-web-metadatos-ocsort}

El pipeline de metadatos original
solo se había ejecutado sobre la configuración óptima de ByteTrack. Para
que el selector de algoritmo de la aplicación web tuviera sentido también
en el explorador de secuencias, \texttt{src/web\_export/compute\_metadata.py}
reproduce exactamente la misma lógica parametrizada por algoritmo, en lugar de duplicar los cinco
scripts de \texttt{src/metadata/} o modificarlos. Los 4 puntos de calibración de
homografía descritos en la sección~\ref{sec:implementacion-homografia}
se reutilizan sin modificación para OC-SORT, porque dependen únicamente
de la posición fija de la cámara en ese tramo de vídeo, no del algoritmo
de tracking.

Para no mezclar ni sobrescribir los artefactos de ByteTrack ya generados, la
salida de OC-SORT se escribe en un árbol de carpetas paralelo,
 con el mismo formato de nombre de fichero que la
rama de ByteTrack. La ejecución sobre las tres secuencias de demostración
produjo por ejemplo porcentajes de posesión por
equipo muy próximos a los ya obtenidos con ByteTrack en el
capítulo~\ref{cap:pruebas} para las mismas secuencias: por ejemplo, en
SNMOT-116 el equipo dominante se queda con el 86,2\,\% de la posesión con
OC-SORT frente al 86,0\,\% con ByteTrack (aunque a qué equipo, 0 o 1, le
corresponde ese porcentaje mayoritario puede no coincidir entre ambas
ejecuciones, ya que la etiqueta numérica que K-means asigna a cada grupo
es arbitraria y no se fija entre ejecuciones independientes). Esta
cercanía es una comprobación indirecta de que la heurística de posesión
por proximidad (sección~\ref{subsec:impl-posesion}) es razonablemente
estable frente al algoritmo de tracking subyacente, ya que ambos parten de
las mismas detecciones y solo difieren en cómo las asocian entre frames.

\subsection{Empaquetado de datos estáticos para la web}
\label{subsec:impl-web-bundle}

\texttt{src/web\_export/export\_bundle.py} genera, para cada una de las
combinaciones de secuencia y algoritmo, el paquete de artefactos que
consume la interfaz, bajo \texttt{frontend/public/data/}:

\begin{itemize}
\item \textbf{Vídeo anotado}: Este vídeo se dibuja
  directamente sobre el fichero de resultados ya calculado con la
  configuración óptima, sin recomputar nada: las cajas se colorean por
  equipo a partir del CSV de clasificación correspondiente, con el balón
  resaltado en un color distinto, sobre el tramo de frames ya usado en la
  homografía o un tramo por
  defecto de longitud comparable. Cada vídeo se reescala a 960\,px de ancho
  antes de codificarse, para mantener un tamaño razonable sin perder
  legibilidad de las cajas.
\item \textbf{Imágenes}: copia directa de los mapas de calor y del mapa
  homografiado ya generados en las secciones~\ref{subsec:impl-web-metadatos-ocsort}
  y~\ref{sec:implementacion-metadatos}.
\item \textbf{\texttt{sequences.json}}: un registro por combinación de
  secuencia y algoritmo, con el recuento de jugadores por equipo, el
  porcentaje de posesión, el top-5 de distancia/velocidad y las rutas a
  vídeo e imágenes; los campos correspondientes a datos que no existen
  para esa combinación concreta se serializan como \texttt{null} en lugar de omitirse, para
  que la interfaz pueda distinguir explícitamente \enquote{no calculado} de
  \enquote{campo ausente por error}.
\end{itemize}

La primera versión de los vídeos, codificada con el \emph{fourcc}
\texttt{mp4v} (MPEG-4 Part 2) de OpenCV, no se reproducía en ningún
navegador (problema P17, Tabla~\ref{tab:impl-problemas}): los navegadores
modernos solo soportan H.264/H.265 dentro de un contenedor
\texttt{.mp4}, no MPEG-4 Part 2. Cambiar el \emph{fourcc} a
\texttt{avc1} ( fue posible porque la instalación de OpenCV usada en este
proyecto está compilada con soporte de FFmpeg) resolvió el problema.

\subsection{Frontend}
\label{subsec:impl-web-frontend}

La parte del frontend es un proyecto independiente creado con
\texttt{vite}, y en producción se compila a HTML/CSS/JS estático.
La configuración de Vite fija \texttt{base:
'./'} para que las rutas del build resultante sean relativas y el sitio
se pueda desplegar en cualquier subruta o servirse en local sin tocar
configuración.

La interfaz reproduce el diseño de la
sección~\ref{subsec:diseno-web-ui}: un componente \texttt{App.tsx} con
una cabecera fija y dos pestañas que
alternan entre \texttt{ComparativaAlgoritmos} y
\texttt{ExploradorSecuencias}. Este último mantiene en estado local la
secuencia y el algoritmo seleccionados, y renderiza cinco componentes de
panel (\texttt{VideoPanel}, \texttt{TeamsPanel}, \texttt{PossessionPanel},
\texttt{HeatmapPanel}, \texttt{HomografiaPanel}) a partir de la entrada de
\texttt{sequences.json} que coincide con esa combinación; ninguno de
estos componentes calcula nada, se limitan a formatear y mostrar los
campos ya presentes en el JSON, incluyendo el caso en que un campo llega
a \texttt{null}, donde el panel muestra una nota explícita de
\enquote{no disponible para esta combinación} para que se lea como una
característica de los datos y no como un fallo de carga.

La paleta de la interfaz sigue un esquema oscuro, con un único color de
acento reutilizado en el subrayado de la pestaña
activa, los selectores de secuencia y algoritmo, y el propio logotipo del
proyecto, el mismo tono ya usado dentro de la tabla comparativa para
resaltar el mejor valor de cada métrica entre los dos algoritmos: una
única decisión de color reutilizada en todo el sitio en lugar de
introducir varias.

\section{Enfoques explorados y descartados}
\label{sec:implementacion-descartados}

Como se ha ido señalando a lo largo de este capítulo y del
proyecto en general, no todo lo probado durante este proyecto
llegó a incorporarse al sistema final. Esta sección reúne, a modo de
resumen, los enfoques explorados y descartados y la
razón de su descarte.

\subsection{Detección del balón con YOLO}
\label{subsec:impl-yolo}

Como se explicó en la sección~\ref{subsec:fundamentos-balon}, la primera
aproximación para separar el balón de los jugadores no fue el filtro
geométrico finalmente adoptado
(sección~\ref{subsec:impl-balon}), sino un detector de objetos entrenado
con aprendizaje profundo, \textbf{YOLOv8n} \citep{redmon2016yolo}, aplicado
directamente sobre los frames.
Se descartó tras las primeras pruebas porque, sin
reentrenamiento específico para balones de fútbol en las condiciones de
cámara e iluminación de SoccerNet, el modelo preentrenado genérico obtenía
resultados muy pobres, con falsos negativos frecuentes y falsos positivos
sobre otros objetos redondeados y claros del campo.

\subsection{Clasificación de equipos con un tercer grupo para el árbitro}

Como se documentó en la sección~\ref{sec:fundamentos-kmeans}, la
clasificación de equipos por K-means asume $k=2$ grupos, lo que
provoca que el árbitro (que tiene una camiseta de color distinto a la de ambos
equipos) se asigne incorrectamente al centroide de color más próximo. Se
valoró una variante con $k=3$, combinada con una heurística adicional que
etiquetaría como árbitro al grupo con menos detecciones totales (teniendo en cuenta
que
solo hay un árbitro principal frente a unos once jugadores por equipo),
pero no llegó a implementarse porque el enfoque con $k=2$ ya ofrecía resultados
suficientemente buenos para los objetivos de este proyecto, y la mejora
quedó pospuesta como línea de trabajo futuro
(capítulo~\ref{cap:conclusiones}) en lugar de invertir tiempo adicional en
una heurística no estrictamente necesaria en esta fase.

\subsection{Homografía generalizada a toda la secuencia}

La homografía real (sección~\ref{sec:implementacion-homografia}) se ha
limitado deliberadamente, desde el principio, a los tramos concretos de
cámara estable identificados en cada secuencia de demostración, y nunca se
intentó generalizarla al resto del dataset. La razón es estructural, no de
tiempo disponible: la mayoría de las secuencias de SoccerNet tienen
movimiento de cámara continuo a lo largo de todo su vídeo,
lo que exigiría recalcular la matriz de homografía de forma dinámica,
frame a frame ( lo que es un problema de calibración de cámara automática
considerablemente más complejo que el cálculo manual y estático empleado
aquí, y que queda fuera del alcance de este proyecto). Acotar expresamente el
alcance de esta técnica a una prueba de concepto, en lugar de forzar su
aplicación al resto del dataset con resultados poco fiables, es coherente
con el riesgo identificado y mitigado en la
Tabla~\ref{tab:gestion-riesgos-tabla} del capítulo~\ref{cap:gestion}.
